\section{Introduction}
\label{sec:intro}

\subsection{Motivation: Scaling Down When Silicon Will Not Scale Up}

The last decade of machine learning has been characterized by a
clear and consistent assumption: that the appropriate response to a
hard problem is a larger model trained on a larger accelerator with
a larger memory budget. This trajectory has produced spectacular
results in language, vision, and multimodal
generation~\cite{patterson2021carbon}. It has also imposed energy,
carbon, and supply-chain costs that are now visible at a
macroeconomic scale: a multi-year global shortage of leading-edge
semiconductors disrupted the automotive, consumer electronics, and
medical-device industries between 2020 and
2024~\cite{burkacky2022chipshortage,khan2021chipshortage}, and the
energy footprint of always-online inference is large and growing
faster than supply~\cite{strubell2019energy,patterson2021carbon}.

Against this backdrop, a complementary research agenda has gained
new urgency. Rather than asking what new accelerator a model
requires, this agenda asks what models can run on the silicon that
is already in mass production, already shipping in tens of billions
of units per year, and already affordable in single-digit-dollar
unit cost: the 8-bit and 16-bit microcontrollers (MCUs) that
populate wearables, sensors, smart appliances, automotive
sub-systems, and industrial telemetry endpoints. This is the
``tinyML'' regime~\cite{banbury2021mlperf}, and it is
infrastructure-agnostic in a way that GPU-scale inference is not.

\subsection{The Bare-Metal MCU Class}

Most published tinyML deployments target ARM Cortex-M class devices
that expose hundreds of kilobytes of SRAM, hardware floating-point
units, single-cycle multipliers, and vendor-optimized neural-network
kernels~\cite{lai2018cmsisnn,david2021tflitemicro,lin2020mcunet}.
These platforms are inexpensive but still relatively powerful.
At the other end of the deployable-silicon spectrum sit the
\emph{bare-metal} MCUs studied in this paper:
\begin{itemize}
\item the 8-bit \arduino{} (ATmega328P), with 32~KB Flash, 2~KB
SRAM, a hardware $8{\times}8$ multiplier, and no floating-point unit;
\item the 16-bit \msp{}, with 16~KB Flash, 512~B SRAM, and
\emph{no hardware multiplier of any kind}. Every multiplication
on this device is software-emulated.
\end{itemize}
These two parts are representative of the low end of the available
silicon supply and remain in active production at unit costs an
order of magnitude below any ARM Cortex-M target. If a recurrent
network can be made to run in real time on the \msp{}, then the same
network can be made to run on essentially any commodity
microcontroller in production today.

\subsection{Why Recurrent Networks, and Why \fastgrnn{}}

Human activity recognition (HAR) from a wrist- or waist-mounted
inertial sensor is a canonical streaming-classification problem
with low input dimensionality (three accelerometer channels), low
sampling rate (50~Hz), and a small number of output classes (six)
\cite{anguita2013public,reyes2015transition}. It is therefore an
ideal benchmark for the bare-metal MCU class: simple enough to fit,
but temporal enough that a recurrent treatment outperforms a
windowed feedforward baseline.

\fastgrnn{}~\cite{kusupati2018fastgrnn} is a gated recurrent cell
designed explicitly for resource-constrained inference. It combines
three orthogonal compression mechanisms -- low-rank weight
factorization, iterative hard-thresholding (IHT)
sparsity~\cite{blumensath2009iht}, and Q15 fixed-point
quantization~\cite{jacob2018integer,han2016deepcompression} -- and
has been reported to match LSTM~\cite{hochreiter1997lstm} accuracy
at a fraction of the parameter count. The original paper, however,
does not provide a public bare-metal reference implementation on a
multiplier-less target, nor does it characterize streaming-mode
behavior.

\subsection{Contributions}

This paper delivers an end-to-end open-source reproduction of
\fastgrnn{} for HAR on bare-metal MCUs and reports four
contributions beyond the original paper:

\begin{itemize}
\item \textbf{Cross-platform deterministic inference.} A single
portable C source compiles on both the 8-bit \arduino{} and the
16-bit \msp{}, producing \emph{bit-equivalent} hidden-state
trajectories and 100\% prediction agreement with a PyTorch
reference across 3{,}399 test windows (seed~0 deployed;
99.91--100\% across five seeds).

\item \textbf{A deployable look-up-table recipe for multiplier-less
targets.} A 256-entry sigmoid/tanh look-up table over the input
range $[-8, +8]$ accelerates full-window inference on the \msp{} by
\textbf{30.5$\times$} (54~s $\rightarrow$ 1.8~s), turning what was
a non-real-time deployment into a comfortable 50~Hz streaming
target with 31\% budget headroom. The recipe is independent of
\fastgrnn{} and applies to any recurrent cell that relies on
$\sigma$ or $\tanh$ activations.

\item \textbf{Characterization of recurrent warm-up latency.} We
quantify, over 100 random test windows, that prediction stability
requires a \textbf{median of 74 samples (1.48~s)} of recurrent
hidden-state evolution, with a worst-case of 125 samples (2.50~s),
regardless of deployment target. This is a property of the
recurrent dynamics, not the underlying hardware, and bounds the
end-to-end user-facing response time of any \fastgrnn{}-class HAR
system in a way that is not discussed in the original paper.

\item \textbf{Hardware energy characterization.} Using a hardware
current sensor (INA226) on the \msp{} target rail, we measure
\textbf{17.7~mW} active inference power, \textbf{$<$0.09~mW} in
LPM3 sleep, and \textbf{31.5~mJ} per 128-sample inference window
with the LUT enabled. Without the LUT, the same window consumes
954~mJ --- a \textbf{96.7\% energy reduction} attributable
entirely to the look-up table's 30.5$\times$ latency gain at
fixed clock frequency. This closes the loop between the model's
algorithmic efficiency and its real-world power envelope, a
characterization absent from the original \fastgrnn{} paper.
\end{itemize}

\noindent\textbf{Code and reproducibility.} All source code,
trained models, exported Q15 weights, per-experiment logs, and
deployment binaries are publicly available
at~\url{https://github.com/emre1998/fastgrnn-har} under the
Apache License 2.0. The complete reproduction takes approximately
two hours of CPU training plus five minutes of MCU deployment.

\subsection{Roadmap}

Section~\ref{sec:related} surveys related compact-RNN and edge-ML
work. Section~\ref{sec:method} formalizes the \fastgrnn{} cell and
the low-rank, sparse, quantized (L-S-Q) compression pipeline.
Section~\ref{sec:setup} describes the experimental setup;
Section~\ref{sec:results} reports accuracy, deployment footprint,
real-time streaming performance, and hardware energy characterization.
Section~\ref{sec:discussion} analyzes the warm-up phenomenon,
discusses cross-platform determinism, and lists limitations.
Section~\ref{sec:conclusion} concludes.
